% Options for packages loaded elsewhere
\PassOptionsToPackage{unicode}{hyperref}
\PassOptionsToPackage{hyphens}{url}
\PassOptionsToPackage{dvipsnames,svgnames,x11names}{xcolor}
\documentclass[
  a4paper,
]{ltjsarticle}
\usepackage{xcolor}
\usepackage[margin=20mm]{geometry}
\usepackage{amsmath,amssymb}
\setcounter{secnumdepth}{-\maxdimen} % remove section numbering
\usepackage{iftex}
\ifPDFTeX
  \usepackage[T1]{fontenc}
  \usepackage[utf8]{inputenc}
  \usepackage{textcomp} % provide euro and other symbols
\else % if luatex or xetex
  \usepackage{unicode-math} % this also loads fontspec
  \defaultfontfeatures{Scale=MatchLowercase}
  \defaultfontfeatures[\rmfamily]{Ligatures=TeX,Scale=1}
\fi
\usepackage{lmodern}
\ifPDFTeX\else
  % xetex/luatex font selection
  \setmainfont[]{DejaVu Serif}
  \setsansfont[]{DejaVu Sans}
  \setmonofont[]{DejaVu Sans Mono}
\fi
% Use upquote if available, for straight quotes in verbatim environments
\IfFileExists{upquote.sty}{\usepackage{upquote}}{}
\IfFileExists{microtype.sty}{% use microtype if available
  \usepackage[]{microtype}
  \UseMicrotypeSet[protrusion]{basicmath} % disable protrusion for tt fonts
}{}
\makeatletter
\@ifundefined{KOMAClassName}{% if non-KOMA class
  \IfFileExists{parskip.sty}{%
    \usepackage{parskip}
  }{% else
    \setlength{\parindent}{0pt}
    \setlength{\parskip}{6pt plus 2pt minus 1pt}}
}{% if KOMA class
  \KOMAoptions{parskip=half}}
\makeatother
\usepackage{longtable,booktabs,array}
\newcounter{none} % for unnumbered tables
\usepackage{calc} % for calculating minipage widths
% Correct order of tables after \paragraph or \subparagraph
\usepackage{etoolbox}
\makeatletter
\patchcmd\longtable{\par}{\if@noskipsec\mbox{}\fi\par}{}{}
\makeatother
% Allow footnotes in longtable head/foot
\IfFileExists{footnotehyper.sty}{\usepackage{footnotehyper}}{\usepackage{footnote}}
\makesavenoteenv{longtable}
\usepackage{graphicx}
\makeatletter
\newsavebox\pandoc@box
\newcommand*\pandocbounded[1]{% scales image to fit in text height/width
  \sbox\pandoc@box{#1}%
  \Gscale@div\@tempa{\textheight}{\dimexpr\ht\pandoc@box+\dp\pandoc@box\relax}%
  \Gscale@div\@tempb{\linewidth}{\wd\pandoc@box}%
  \ifdim\@tempb\p@<\@tempa\p@\let\@tempa\@tempb\fi% select the smaller of both
  \ifdim\@tempa\p@<\p@\scalebox{\@tempa}{\usebox\pandoc@box}%
  \else\usebox{\pandoc@box}%
  \fi%
}
% Set default figure placement to htbp
\def\fps@figure{htbp}
\makeatother
\setlength{\emergencystretch}{3em} % prevent overfull lines
\providecommand{\tightlist}{%
  \setlength{\itemsep}{0pt}\setlength{\parskip}{0pt}}
\usepackage{xurl}
\usepackage{float}
\usepackage{bookmark}
\IfFileExists{xurl.sty}{\usepackage{xurl}
\usepackage{float}}{} % add URL line breaks if available
\urlstyle{same}
\hypersetup{
  pdftitle={三次元正方配列の三角形ループ積における電弱対称性破れとパリティ最大破れ（Part IV）},
  pdfauthor={木原 範昭（株式会社WFシステム）　ORCID 0009-0004-6753-4020},
  colorlinks=true,
  linkcolor={Maroon},
  filecolor={Maroon},
  citecolor={Blue},
  urlcolor={Blue},
  pdfcreator={LaTeX via pandoc}}

\title{三次元正方配列の三角形ループ積における電弱対称性破れとパリティ最大破れ（Part
IV）}
\usepackage{etoolbox}
\makeatletter
\providecommand{\subtitle}[1]{% add subtitle to \maketitle
  \apptocmd{\@title}{\par {\large #1 \par}}{}{}
}
\makeatother
\subtitle{― 八元数 \(G_2 \supset SO(4)\) による 100\%:0\%
カイラル分離、\(\mathrm{Tri}(T)\)
動的ヒッグス質量生成、\(W^\pm/Z^0/\gamma\)
スペクトル、および素のワインバーグ角
\(\sin^2\theta_W^{\mathrm{bare}} = 3/8\) ―}
\author{木原 範昭（株式会社WFシステム）　ORCID 0009-0004-6753-4020}
\date{2026年11月28日　Version DOI 10.5281/zenodo.24626593}

\begin{document}
\maketitle

{
\setcounter{tocdepth}{2}
\tableofcontents
}
\textbf{著者：} 木原 範昭（株式会社WFシステム）\\
\textbf{ORCID：} 0009-0004-6753-4020\\
\textbf{日付：} 2026年11月28日\\
\textbf{Version DOI：} 10.5281/zenodo.24626593\\
\textbf{Concept DOI：} 10.5281/zenodo.24626592\\
\textbf{位置づけ：} 三角形ループ積シリーズ Part I（Concept DOI:
\texttt{10.5281/zenodo.24622200}）・Part II（Concept DOI:
\texttt{10.5281/zenodo.24622565}）・Part III（Concept DOI:
\texttt{10.5281/zenodo.24625962}）の続編となる\textbf{電弱対称性破れ・パリティ最大破れ検証論文（Part
IV）}。Part III で特定した外側 \(O_8^{\mathrm{ext}}\) の \texttt{1+3+4}
分割（純粋 4 空間
\(SU(2)\)）を起点に、標準模型の電弱セクター（\(SU(2)_L \times U(1)_Y \to U(1)_{\mathrm{em}}\)、パリティ最大破れ、\(W/Z\)
ボソン質量獲得、ワインバーグ角）を乱数なしの決定論的計算（\texttt{tri\_loop\_part4\_electroweak\_parity\_v0.py}）で直接検証する。

\begin{center}\rule{0.5\linewidth}{0.5pt}\end{center}

\subsection{要旨}\label{ux8981ux65e8}

Part I--III {[}P1--P3{]} により、背景時空や作用を仮定しない実数 3
階テンソルの三角形ループ積
\(\mathrm{Tri}(T)_{ijk} = \sum_{p,q,r} T_{ipq} T_{qjr} T_{rpk}\)
から、重力（\texttt{1+3}
過去光円錐と潮汐）・強い力（\(SU(3)_C\)）・電磁気（\(U(1)_Q\)）・1 世代
CAR フェルミオン・分解能ゲージが統一接続され、さらに Part III で外側
\(O_8^{\mathrm{ext}}\) の余り 4 空間 \(\{e_4,e_5,e_6,e_7\}\) に 3
空間不変な純粋 3 次元 \(SU(2)\)
が存在することが特定された。しかし標準模型 1 世代を閉じるためには、(1)
なぜ弱い力 \(SU(2)_L\)
だけが片方のカイラリティにのみ結合するのか（パリティ最大破れ）、(2)
どのようにハイパーチャージ \(U(1)_Y\)
が全ゲージ群と可換に組まれるか、(3) 実際の \(\mathrm{Tri}(T)\)
動力学でヒッグス凝縮により \(W/Z\) 質量と無質量光子 \(\gamma\)
が分離するか、の 3 点が未解決であった。

本 Part IV
では、乱数なしの決定論的検証（\texttt{tri\_loop\_part4\_electroweak\_parity\_v0.py}）により以下
5 点を確立・報告する。

\begin{enumerate}
\def\labelenumi{\arabic{enumi}.}
\tightlist
\item
  \textbf{パリティ最大破れの厳密な代数的起源（\(G_2\) カイラル 100\%:0\%
  定理）}：4 空間 \(\{e_4,e_5,e_6,e_7\}\) の回転群
  \(SO(4) \cong SU(2)_{\mathrm{SD}} \times SU(2)_{\mathrm{ASD}}\)（自己双対＝+カイラル、反自己双対＝−カイラル）に対し、八元数の自己同型群
  \(G_2\)（14 次元）の微分制約を課すと、3 空間 \(\{e_1,e_2,e_3\}\)
  を完全に固定する内部 \(SU(2)_{\mathrm{pure4}}\) は
  \textbf{自己双対（SD）占有率
  \texttt{1.000000000000}・反自己双対（ASD）占有率
  \texttt{0.000000000000}} に厳密一致する（\(G_2\)
  微分残差・リー代数残差とも
  \textbf{\texttt{0.000e+00}}）。逆カイラル（ASD）側は \texttt{100\%}
  空間回転 \(SO(3)\) にロックされ独立内部ゲージになれない。
\item
  \textbf{標準模型全 12 ゲージ生成子
  \(SU(3)_C \times SU(2)_L \times U(1)_Y\)
  の完全閉包とパリティ選択則}：厳密可換ハイパーチャージ
  \(Y = I_8 \otimes Q_{\mathrm{int}} - \frac{1}{2} P_4 \otimes I_8\)（ゲルマン＝西島
  \(Q_{\mathrm{em}} = T_3^L + Y\)）を定義すると、\([SO(3)_{\mathrm{ext}}, \mathrm{SM}_{12}]\)（全
  36 組）・\([SU(3)_C, SU(2)_L]\)（24 組）・\([SU(3)_C, U(1)_Y]\)（8
  組）・\([SU(2)_L, U(1)_Y]\)（3 組）のすべてが最大誤差
  \textbf{\texttt{0.000e+00}} で可換となる。\(W_1\)
  回転（\(\theta=0.65\)）では左巻きダブレット（\(e_4\) 側）が
  \texttt{0.323572} で強く結合する一方、右巻きシングレット（\(e_0, e_1\)
  側）は \textbf{\texttt{0.000e+00}}（厳密非結合）となる。
\item
  \textbf{\(\mathrm{Tri}(T)\) 実動力学によるヒッグス凝縮と \(W\)
  ボソン質量生成 vs 残余無質量}：真空 \(H = O_8 \otimes M_8\) では
  \(SU(2)_L\)・残余 8D 対称群とも対称性破れノルムが
  \(\le 3.8 \times 10^{-14}\)（完全不変・無質量）である。ここに 4
  空間サイト \(A=4\) 局在のヒッグス凝縮体
  P4\(_4\)（\(\epsilon_H = 0.05 \to 0.22, \tau=40\)）を
  \(\mathrm{Tri}(T)\) で発展させると、\(W_1, W_2, W_3\) の破れノルム
  \(\|W_a(U)\|_F\) は
  \textbf{\texttt{0.012593\ \$\textbackslash{}to\$\ 0.025105\ \$\textbackslash{}to\$\ 0.035323\ \$\textbackslash{}to\$\ 0.035517}}
  へと 3 成分 6 桁完全縮退で立ち上がる一方、\(e_4\) 固定の残余 8D
  対称群は全振幅で \textbf{\texttt{3.82e-14}}（厳密無質量）を維持する。
\item
  \textbf{素のワインバーグ角 \(\sin^2\theta_W^{\mathrm{bare}} = 3/8\) と
  4×4 電弱質量スペクトル（\(\gamma, W^\pm, Z^0\)）}：\(SU(5)/SO(10)\)
  を仮定せず 64 次元状態空間上の素のトレース比だけで
  \(\mathrm{Tr}(T_3^2) = 8\)、\(\mathrm{Tr}(Y^2) = 40/3\)、\(\mathrm{Tr}(Q_{\mathrm{em}}^2) = 64/3\)
  となり、\textbf{\(\sin^2\theta_W^{\mathrm{bare}} = \mathrm{Tr}(T_3^2)/\mathrm{Tr}(Q_{\mathrm{em}}^2) = 3/8 = 0.375000\)}（大統一理論の標準理論値
  \(3/8\) と 6 桁一致）が自動導出される。この幾何比
  \(g'/g = \sqrt{3/5}\)（調整パラメータゼロ）を複素ヒッグス二重項（\(T_3^L=-1/2, Y=+1/2 \Rightarrow Q_{\mathrm{em}}=0\)）の
  4×4 質量行列 \(M^2\) に適用すると、固有値
  \([m_\gamma^2, m_{W^+}^2, m_{W^-}^2, m_Z^2] = [\mathbf{-2.082\times 10^{-17}},\; \mathbf{0.250000},\; \mathbf{0.250000},\; \mathbf{0.400000}]\)、\textbf{\(m_W^2 / m_Z^2 = 5/8 = 0.625000 = \cos^2\theta_W^{\mathrm{bare}}\)}
  が厳密に得られる。
\item
  \textbf{内側時間結晶反転（\(M_8\) vs
  \(M_8^T\)）の誠実監査}：背景と凝縮体を同時に鏡映反転すると
  \(\|P\| = 0.046825\)・\(\|W_1(U)\|_F = 0.036160\) が厳密一致し、Part I
  §8.3 の \(-1.41\%\)
  分裂が背景と粒子の\textbf{相対カイラリティ差}に由来する健全な対称性であることを確認した。
\end{enumerate}

\begin{center}\rule{0.5\linewidth}{0.5pt}\end{center}

\subsection{1. 導入：標準模型 1
世代を閉じるための電弱・パリティ課題}\label{ux5c0eux5165ux6a19ux6e96ux6a21ux578b-1-ux4e16ux4ee3ux3092ux9589ux3058ux308bux305fux3081ux306eux96fbux5f31ux30d1ux30eaux30c6ux30a3ux8ab2ux984c}

Part II {[}P2{]} は
\(O_8^{\mathrm{ext}} \otimes O_8^{\mathrm{int}}\)（64 実自由度）から
\(SU(3)_C \times U(1)_Q\) と 1 世代 CAR 電荷 \texttt{1:3:3:1}
を導き、Part III {[}P3{]} は外側 \(O_8^{\mathrm{ext}}\) の
\texttt{1+3+4} 分割から 3 次元空間と可換な純粋 3 次元 \(SU(2)\)
を特定した。ここで標準模型のゲージ群
\(SU(3)_C \times SU(2)_L \times U(1)_Y\)
を完全に閉じるために残っていた核心の問いは 3 つである：

\begin{enumerate}
\def\labelenumi{\arabic{enumi}.}
\tightlist
\item
  \textbf{なぜパリティが最大に破れるのか？}：4 空間
  \(\{e_4,e_5,e_6,e_7\}\) の回転群
  \(SO(4) \cong SU(2)_L \times SU(2)_R\) には本来左右 2 つの \(SU(2)\)
  がある。なぜその片方だけが内部の弱い力になり、反対側が消えるのか。
\item
  \textbf{ハイパーチャージ \(U(1)_Y\)
  とワインバーグ角は調整なしで決まるか？}：\(SU(3)_C \times SU(2)_L\)
  と厳密可換な \(U(1)_Y\) を構成したとき、混合角 \(\sin^2\theta_W\)
  の幾何学的基準値はいくつになるか。
\item
  \textbf{\(\mathrm{Tri}(T)\) の実動力学で自発的対称性の破れと
  \(W/Z/\gamma\)
  質量分離が起きるか？}：ポテンシャルを手で置かずに、局在凝縮体によって
  \(SU(2)_L\) だけが破れて残余対称が無質量に残るか。
\end{enumerate}

本論文でも Part II・III と同じく、各主張の資格を
\textbf{【既知】・【Part I--III 確立】・【Part IV
検証】・【仮説・未検証】} の 4 段階で明示する。

\begin{center}\rule{0.5\linewidth}{0.5pt}\end{center}

\subsection{\texorpdfstring{2. \(G_2 \supset SO(4)\) による 100\%:0\%
カイラル分離（パリティ最大破れの代数的起源）【Part IV 検証
EW-1】}{2. G\_2 \textbackslash supset SO(4) による 100\%:0\% カイラル分離（パリティ最大破れの代数的起源）【Part IV 検証 EW-1】}}\label{g_2-supset-so4-ux306bux3088ux308b-1000-ux30abux30a4ux30e9ux30ebux5206ux96e2ux30d1ux30eaux30c6ux30a3ux6700ux5927ux7834ux308cux306eux4ee3ux6570ux7684ux8d77ux6e90part-iv-ux691cux8a3c-ew-1}

4 次元実空間 \(\{e_4, e_5, e_6, e_7\}\) の回転代数
\(\mathfrak{so}(4) \cong \mathfrak{su}(2)_{\mathrm{SD}} \oplus \mathfrak{su}(2)_{\mathrm{ASD}}\)
は、3 次元の自己双対（SD：+カイラル）部分代数と 3
次元の反自己双対（ASD：−カイラル）部分代数に直交分解される【既知】。

Part III で抽出した \(G_2\) 内の \texttt{3+4} 保存部分群
\(SO(4) = SU(2)_{\mathrm{pure4}} \times SU(2)_{\mathrm{space}}\)（6
次元）の 4×4 ブロックを、SD と ASD の両部分空間へ射影した結果（図
1(a)）：

{\def\LTcaptype{none} % do not increment counter
\begin{longtable}[]{@{}
  >{\raggedright\arraybackslash}p{(\linewidth - 6\tabcolsep) * \real{0.2222}}
  >{\raggedright\arraybackslash}p{(\linewidth - 6\tabcolsep) * \real{0.2222}}
  >{\centering\arraybackslash}p{(\linewidth - 6\tabcolsep) * \real{0.2778}}
  >{\centering\arraybackslash}p{(\linewidth - 6\tabcolsep) * \real{0.2778}}@{}}
\toprule\noalign{}
\begin{minipage}[b]{\linewidth}\raggedright
\(G_2\) 内の 3 次元部分群
\end{minipage} & \begin{minipage}[b]{\linewidth}\raggedright
3 空間 \(\{e_1,e_2,e_3\}\) への作用
\end{minipage} & \begin{minipage}[b]{\linewidth}\centering
4 空間 SD（+カイラル）占有率
\end{minipage} & \begin{minipage}[b]{\linewidth}\centering
4 空間 ASD（−カイラル）占有率
\end{minipage} \\
\midrule\noalign{}
\endhead
\bottomrule\noalign{}
\endlastfoot
\textbf{\(SU(2)_{\mathrm{pure4}}\)（内部弱ゲージ）} & 厳密ゼロ（最大
\texttt{8.327e-17}） & \textbf{\texttt{1.000000000000}} &
\textbf{\texttt{0.000000000000}} \\
\textbf{\(SU(2)_{\mathrm{space}}\)（空間回転随伴）} & \(SO(3)\) 空間回転
& \textbf{\texttt{0.000000000000}} & \textbf{\texttt{1.000000000000}} \\
\end{longtable}
}

\begin{itemize}
\tightlist
\item
  標準 \(SU(2)_L\) 基底 \((T_1, T_2, T_3)\) の \(G_2\) 微分残差は
  \textbf{\texttt{0.000e+00}}、リー代数残差も
  \textbf{\texttt{0.000e+00}}。
\item
  \textbf{物理的帰結}：八元数の乗算表（\(G_2\) 自己同型）は、3
  次元空間と可換な独立内部ゲージ \(SU(2)\) として
  \textbf{片側のカイラリティ（SD）だけ}を許し、逆側のカイラリティ（ASD）を
  100\% 外部 3 次元空間回転 \(SO(3)\)
  へロックする。これが、弱い力が左巻き成分のみに結合し右巻き成分に結合しない（パリティ最大破れ）ことの厳密な代数的起源である【Part
  IV 検証】＋【仮説】。
\end{itemize}

\begin{center}\rule{0.5\linewidth}{0.5pt}\end{center}

\subsection{3. 標準模型全 12
ゲージ生成子の厳密可換性とパリティ選択則【Part IV 検証
EW-2】}\label{ux6a19ux6e96ux6a21ux578bux5168-12-ux30b2ux30fcux30b8ux751fux6210ux5b50ux306eux53b3ux5bc6ux53efux63dbux6027ux3068ux30d1ux30eaux30c6ux30a3ux9078ux629eux5247part-iv-ux691cux8a3c-ew-2}

64 次元空間 \(O_8^{\mathrm{ext}} \otimes O_8^{\mathrm{int}}\) 上で、内側
\(SU(3)_C\)（8 生成子）・外側 \(SU(2)_L\)（3 生成子）に加え、4 空間射影
\(P_4 = \mathrm{diag}(0,0,0,0,1,1,1,1)\) を用いたハイパーチャージ
\[Y \equiv I_8 \otimes Q_{\mathrm{int}} - \frac{1}{2}\, P_4 \otimes I_8, \qquad Q_{\mathrm{em}} = T_3^L + Y \tag{1}\]
を構成した。全ゲージ群および外部空間回転 \(SO(3)_{\mathrm{ext}}\)
間の交換子最大誤差は以下の通りである：

\begin{itemize}
\tightlist
\item
  \([SO(3)_{\mathrm{ext}},\; SU(3)_C \times SU(2)_L \times U(1)_Y]\)（全
  \(3 \times 12 = 36\) 組）：\textbf{\texttt{0.000e+00}}
\item
  \([SU(3)_C,\; SU(2)_L]\)（全 24 組）：\textbf{\texttt{0.000e+00}}
\item
  \([SU(3)_C,\; U(1)_Y]\)（全 8 組）：\textbf{\texttt{0.000e+00}}
\item
  \([SU(2)_L,\; U(1)_Y]\)（全 3 組）：\textbf{\texttt{0.000e+00}}
\end{itemize}

さらに \(W_1\) 有限回転 \(U_W = \exp(0.65\, T_1^L \otimes I_8)\)
を作用させると、外側 4 空間の左巻きダブレット状態（\(e_4\) 側）は
\(\|U_W\psi_L - \psi_L\| = \mathbf{0.323572}\) と強く結合する一方、外側
\texttt{1+3} 空間の右巻きシングレット状態（\(e_0, e_1\) 側）は
\textbf{\texttt{0.000e+00}}（厳密非結合）となる。

\begin{center}\rule{0.5\linewidth}{0.5pt}\end{center}

\subsection{\texorpdfstring{4. \(\mathrm{Tri}(T)\)
実動力学によるヒッグス凝縮・\(W/Z/\gamma\)
質量スペクトル・\(\sin^2\theta_W^{\mathrm{bare}} = 3/8\)}{4. \textbackslash mathrm\{Tri\}(T) 実動力学によるヒッグス凝縮・W/Z/\textbackslash gamma 質量スペクトル・\textbackslash sin\^{}2\textbackslash theta\_W\^{}\{\textbackslash mathrm\{bare\}\} = 3/8}}\label{mathrmtrit-ux5b9fux52d5ux529bux5b66ux306bux3088ux308bux30d2ux30c3ux30b0ux30b9ux51ddux7e2ewzgamma-ux8ceaux91cfux30b9ux30daux30afux30c8ux30ebsin2theta_wmathrmbare-38}

\subsubsection{\texorpdfstring{4.1 \(\mathrm{Tri}(T)\) 実動力学における
\(W_1, W_2, W_3\) 質量生成 vs 残余無質量【Part IV 検証
EW-3(a,b)】}{4.1 \textbackslash mathrm\{Tri\}(T) 実動力学における W\_1, W\_2, W\_3 質量生成 vs 残余無質量【Part IV 検証 EW-3(a,b)】}}\label{mathrmtrit-ux5b9fux52d5ux529bux5b66ux306bux304aux3051ux308b-w_1-w_2-w_3-ux8ceaux91cfux751fux6210-vs-ux6b8bux4f59ux7121ux8ceaux91cfpart-iv-ux691cux8a3c-ew-3ab}

メキシカンハット型ポテンシャルを外から一切置かず、Part I の真空時間結晶
\(H = O_8 \otimes M_8\) に対し、外側 4 空間のサイト \(A = 4\)
に局在凝縮体 P4\(_4\)（振幅 \(\epsilon_H\)）を投入して
\(\mathrm{Tri}(T)\) を \(\tau = 40\) ステップ発展させた安定状態
\(U_{\mathrm{higgs}}\) 上で、各生成子の対称性破れ秩序パラメータ
\(\|G(U_{\mathrm{higgs}})\|_F\)（ヒッグス運動項 \(|D_\mu\Phi|^2\)
に対応）を計測した（表 1、図 1(b)）。

{\def\LTcaptype{none} % do not increment counter
\begin{longtable}[]{@{}
  >{\centering\arraybackslash}p{(\linewidth - 10\tabcolsep) * \real{0.1667}}
  >{\centering\arraybackslash}p{(\linewidth - 10\tabcolsep) * \real{0.1667}}
  >{\centering\arraybackslash}p{(\linewidth - 10\tabcolsep) * \real{0.1667}}
  >{\centering\arraybackslash}p{(\linewidth - 10\tabcolsep) * \real{0.1667}}
  >{\centering\arraybackslash}p{(\linewidth - 10\tabcolsep) * \real{0.1667}}
  >{\centering\arraybackslash}p{(\linewidth - 10\tabcolsep) * \real{0.1667}}@{}}
\toprule\noalign{}
\begin{minipage}[b]{\linewidth}\centering
ヒッグス凝縮振幅 \(\epsilon_H\)
\end{minipage} & \begin{minipage}[b]{\linewidth}\centering
\(\|W_1(U)\|_F\)
\end{minipage} & \begin{minipage}[b]{\linewidth}\centering
\(\|W_2(U)\|_F\)
\end{minipage} & \begin{minipage}[b]{\linewidth}\centering
\(\|W_3(U)\|_F\)
\end{minipage} & \begin{minipage}[b]{\linewidth}\centering
\(e_4\) 固定残余 8D 最大
\end{minipage} & \begin{minipage}[b]{\linewidth}\centering
凝縮エネルギー \(E_{\mathrm{probe}}\)
\end{minipage} \\
\midrule\noalign{}
\endhead
\bottomrule\noalign{}
\endlastfoot
\texttt{0.00}（真空 \(H\)） & \texttt{0.000000} & \texttt{0.000000} &
\texttt{0.000000} & \texttt{3.802e-14} & \texttt{0.000000} \\
\texttt{0.05} & \textbf{\texttt{0.012593}} & \textbf{\texttt{0.012593}}
& \textbf{\texttt{0.012593}} & \texttt{3.821e-14} & \texttt{0.000294} \\
\texttt{0.10} & \textbf{\texttt{0.025105}} & \textbf{\texttt{0.025105}}
& \textbf{\texttt{0.025105}} & \texttt{3.821e-14} & \texttt{0.000964} \\
\texttt{0.16} & \textbf{\texttt{0.035323}} & \textbf{\texttt{0.035323}}
& \textbf{\texttt{0.035323}} & \texttt{3.828e-14} & \texttt{0.001510} \\
\texttt{0.22} & \textbf{\texttt{0.035517}} & \textbf{\texttt{0.035517}}
& \textbf{\texttt{0.035517}} & \texttt{3.821e-14} & \texttt{0.001055} \\
\end{longtable}
}

真空 \(H\) では完全無質量だった \(W_1, W_2, W_3\)
が、\(\mathrm{Tri}(T)\) 動力学による凝縮の成立と同時に 3 成分 6
桁完全縮退で非ゼロの破れノルムを獲得する一方、\(e_4\) を固定する残余 8
次元対称群は全振幅で機械ゼロ（\texttt{3.82e-14}＝無質量）のまま保存される。

\begin{figure}[H]
\centering
\includegraphics[width=0.96\linewidth,height=\textheight,keepaspectratio,alt={Part IV 検証 (1)：(a) G\_2 \textbackslash supset SO(4) における 100\%:0\% カイラル分離（パリティ最大破れの代数的起源）、(b) \textbackslash mathrm\{Tri\}(T) 実動力学におけるヒッグス凝縮 P4\_4 による W\_1,W\_2,W\_3 質量生成 vs 残余 8D 対称群の厳密無質量保存。}]{figs_ja/fig_p4_parity_and_higgs_mass_ja.png}
\caption{Part IV 検証 (1)：(a) \(G_2 \supset SO(4)\) における 100\%:0\%
カイラル分離（パリティ最大破れの代数的起源）、(b) \(\mathrm{Tri}(T)\)
実動力学におけるヒッグス凝縮 P4\(_4\) による \(W_1,W_2,W_3\) 質量生成 vs
残余 8D 対称群の厳密無質量保存。}
\end{figure}

\subsubsection{\texorpdfstring{4.2 素のワインバーグ角
\(\sin^2\theta_W^{\mathrm{bare}} = 3/8\) と 4×4 質量スペクトル【Part IV
検証 EW-3(c),
EW-5】}{4.2 素のワインバーグ角 \textbackslash sin\^{}2\textbackslash theta\_W\^{}\{\textbackslash mathrm\{bare\}\} = 3/8 と 4×4 質量スペクトル【Part IV 検証 EW-3(c), EW-5】}}\label{ux7d20ux306eux30efux30a4ux30f3ux30d0ux30fcux30b0ux89d2-sin2theta_wmathrmbare-38-ux3068-44-ux8ceaux91cfux30b9ux30daux30afux30c8ux30ebpart-iv-ux691cux8a3c-ew-3c-ew-5}

64 次元状態空間 \(O_8^{\mathrm{ext}} \otimes O_8^{\mathrm{int}}\) 上で式
(1) のトレース二乗和を計算すると（図 2(b)）：
\[\mathrm{Tr}(T_3^2) = 8.000000,\qquad \mathrm{Tr}(Y^2) = \frac{40}{3} = 13.333333,\qquad \mathrm{Tr}(Q_{\mathrm{em}}^2) = \frac{64}{3} = 21.333333 \tag{2}\]
となり、\(SU(5)\) や \(SO(10)\)
大統一群を一切仮定せずとも、幾何学的正規化比から
\[\sin^2\theta_W^{\mathrm{bare}} = \frac{\mathrm{Tr}(T_3^2)}{\mathrm{Tr}(Q_{\mathrm{em}}^2)} = \frac{8}{64/3} = \frac{3}{8} = \mathbf{0.375000}, \qquad \frac{g'}{g} = \sqrt{\frac{\mathrm{Tr}(T_3^2)}{\mathrm{Tr}(Y^2)}} = \sqrt{\frac{3}{5}} \tag{3}\]
が調整パラメータゼロで導かれる {[}E4{]}。

この幾何比 \(g'/g = \sqrt{3/5}\) を用い、外側 4
空間の複素ヒッグス二重項の中性成分（\(T_3^L = -1/2, Y = +1/2 \Rightarrow Q_{\mathrm{em}} = 0\)）に対する
4×4 電弱質量行列 \(M_{ab}^2 = \langle G_a\phi_H, G_b\phi_H \rangle\)
を対角化すると（図 2(a)）：
\[[m_\gamma^2,\; m_{W^+}^2,\; m_{W^-}^2,\; m_Z^2] = [\mathbf{-2.082\times 10^{-17}},\; \mathbf{0.250000},\; \mathbf{0.250000},\; \mathbf{0.400000}] \tag{4}\]
\[\frac{m_W^2}{m_Z^2} = \mathbf{0.625000} = \frac{5}{8} = \cos^2\theta_W^{\mathrm{bare}} \tag{5}\]
となり、厳密無質量な光子 \(\gamma\)・縮退した荷電ボソン
\(W^\pm\)・より重い中性ボソン \(Z^0\) の 4
状態スペクトルが完全再現される。

\begin{figure}[H]
\centering
\includegraphics[width=0.96\linewidth,height=\textheight,keepaspectratio,alt={Part IV 検証 (2)：(a) 複素ヒッグス二重項の 4×4 電弱質量スペクトル（光子厳密無質量、m\_W\^{}2/m\_Z\^{}2 = 5/8）、(b) 64 次元空間上のトレース二乗和から調整なしで導かれる素のワインバーグ角 \textbackslash sin\^{}2\textbackslash theta\_W\^{}\{\textbackslash mathrm\{bare\}\} = 3/8。}]{figs_ja/fig_p4_ew_spectrum_and_weinberg_ja.png}
\caption{Part IV 検証 (2)：(a) 複素ヒッグス二重項の 4×4
電弱質量スペクトル（光子厳密無質量、\(m_W^2/m_Z^2 = 5/8\)）、(b) 64
次元空間上のトレース二乗和から調整なしで導かれる素のワインバーグ角
\(\sin^2\theta_W^{\mathrm{bare}} = 3/8\)。}
\end{figure}

\subsubsection{\texorpdfstring{4.3 内側時間結晶反転（\(M_8\) vs
\(M_8^T\)）の誠実監査【Part IV 検証
EW-4】}{4.3 内側時間結晶反転（M\_8 vs M\_8\^{}T）の誠実監査【Part IV 検証 EW-4】}}\label{ux5185ux5074ux6642ux9593ux7d50ux6676ux53cdux8ee2m_8-vs-m_8tux306eux8aa0ux5b9fux76e3ux67fbpart-iv-ux691cux8a3c-ew-4}

背景 \(H\) と凝縮体 P4\(_4\)
の内側時間結晶を同時に鏡映反転（\(M_8 \to M_8^T\)）して \(\tau = 30\)
発展させたところ、\(\|P\| = 0.046825\) および
\(\|W_1(U)\|_F = 0.036160\) が両者で完全一致した。これは、Part I §8.3 の
\(-1.41\%\)
分裂が「系全体の鏡映」ではなく「背景と粒子の回転向きが逆向き（相対カイラリティ差）」のときにのみ現れる物理的効果であることを確認する結果である。

\begin{center}\rule{0.5\linewidth}{0.5pt}\end{center}

\subsection{5.
残課題の厳密な切り分け（本論文でまだ主張しないこと）}\label{ux6b8bux8ab2ux984cux306eux53b3ux5bc6ux306aux5207ux308aux5206ux3051ux672cux8ad6ux6587ux3067ux307eux3060ux4e3bux5f35ux3057ux306aux3044ux3053ux3068}

本 Part IV により、標準模型 1 世代のゲージ群
\(SU(3)_C \times SU(2)_L \times U(1)_Y\)・パリティ最大破れの代数的起源・\(\mathrm{Tri}(T)\)
動的ヒッグス破れ・素のワインバーグ角 \(3/8\)
が閉じた。以下の課題は【仮説・未検証】として次期論文（Part V
以降）へ切り分ける：

\begin{enumerate}
\def\labelenumi{\arabic{enumi}.}
\tightlist
\item
  \textbf{繰り込み群走行（\(\sin^2\theta_W: 3/8 \to 0.231\)）}：高エネルギー幾何値
  \(3/8\) から低エネルギー実測値 \(\approx 0.231\)
  へのスケール依存性の導出。
\item
  \textbf{多殻大格子 \(O_8^{\otimes K}\) の連続極限（Part V
  予定）}：\(1/r\) 重力・\(1/r^2\) クーロン・光円錐速度
  \(c\)・アインシュタイン方程式・累積重力時間遅延 \(g_{00}(r)\)。
\item
  \textbf{動的散乱と多体量子相関（Part VI 予定）}：2
  粒子の引力・斥力・運動量交換および空間分離 CHSH ベル不等式の動的破れ。
\item
  \textbf{3 世代構造と精密質量比・結合定数（Part VII
  予定）}：例外ジョルダン代数 \(J_3(\mathbb{O})\) 等からの 3
  世代同定、精密質量比、および微細構造定数
  \(\alpha \approx 1/137.036\)・\(\alpha_s\) の導出。
\end{enumerate}

\begin{center}\rule{0.5\linewidth}{0.5pt}\end{center}

\subsection{6.
再現手順と公開ファイル}\label{ux518dux73feux624bux9806ux3068ux516cux958bux30d5ux30a1ux30a4ux30eb}

\begin{itemize}
\tightlist
\item
  検証スクリプト：\texttt{tri\_loop\_part4\_electroweak\_parity\_v0.py}（乱数なし・実行時間約
  15 秒）
\item
  数値ログ：\texttt{results/tri\_loop\_part4\_electroweak\_parity\_v0.txt}
\item
  作図スクリプト：\texttt{plot\_part4\_figures\_v0.py}
\item
  図出力（図 1・図
  2）：\texttt{results/fig\_p4\_parity\_and\_higgs\_mass\_\{ja,en\}.png},
  \texttt{results/fig\_p4\_ew\_spectrum\_and\_weinberg\_\{ja,en\}.png}
\end{itemize}

\begin{center}\rule{0.5\linewidth}{0.5pt}\end{center}

\subsection{参考文献}\label{ux53c2ux8003ux6587ux732e}

\subsubsection{A. 自己引用}\label{a.-ux81eaux5df1ux5f15ux7528}

\begin{itemize}
\tightlist
\item
  \textbf{{[}P1{]}} 木原範昭,
  \emph{三次元正方配列の三角形ループ積からの量子・ゲージ・重力・フェルミオン構造の自発創発（統合完成論文
  Part I）}, Zenodo (2026-11-28), Concept DOI:
  \texttt{10.5281/zenodo.24622200}.
\item
  \textbf{{[}P2{]}} 木原範昭,
  \emph{三次元正方配列の三角形ループ積における重力・電磁気・QCD・ディラックの統一接続仮説（Part
  II）}, Zenodo (2026-11-28), Concept DOI:
  \texttt{10.5281/zenodo.24622565}.
\item
  \textbf{{[}P3{]}} 木原範昭,
  \emph{三次元正方配列の三角形ループ積における時空・内部空間読出し接続仮説（Part
  III）}, Zenodo (2026-11-28), Concept DOI:
  \texttt{10.5281/zenodo.24625962}.
\end{itemize}

\subsubsection{B.
外部一次文献}\label{b.-ux5916ux90e8ux4e00ux6b21ux6587ux732e}

\begin{itemize}
\tightlist
\item
  \textbf{{[}E1{]}} S. Weinberg, \emph{A Model of Leptons}, Phys.
  Rev.~Lett. \textbf{19}, 1264--1266 (1967); A. Salam, in
  \emph{Elementary Particle Theory}, ed.~N. Svartholm (1968).
\item
  \textbf{{[}E2{]}} P. W. Higgs, \emph{Broken Symmetries and the Masses
  of Gauge Bosons}, Phys. Rev.~Lett. \textbf{13}, 508--509 (1964).
\item
  \textbf{{[}E3{]}} C. Furey,
  \emph{\(SU(3)_C \times SU(2)_L \times U(1)_Y (\times U(1)_X)\) as a
  symmetry of division algebraic ladder operators}, Eur. Phys. J. C
  \textbf{78}, 375 (2018), arXiv:1806.00612.
\item
  \textbf{{[}E4{]}} H. Georgi and S. L. Glashow, \emph{Unity of All
  Elementary-Particle Forces}, Phys. Rev.~Lett. \textbf{32}, 438--441
  (1974).
\item
  \textbf{{[}E5{]}} J. C. Baez, \emph{The Octonions}, Bull. Amer. Math.
  Soc. \textbf{39}(2), 145--205 (2002), arXiv:math/0105155.
\end{itemize}

\end{document}
